% 2-4-deep.tex
%  Budget: 2pp.
%  Rules: every numeral in \lr{}; cite only from references.bib;
%  em-dash must be \lr{---} (bare --- is a tofu box in B Nazanin).

\section{یادگیری عمیق}

ورود یادگیری عمیق به این حوزه، هم مجموعه مدل‌ها را گسترش داد و هم پرسش
ارزیابی را حادتر کرد. مدل \lr{LSTM} در پژوهش حاضر نماینده این خانواده
است.

لاگو و همکاران \cite{lago_forecasting_2018} چهار مدل یادگیری عمیق را برای
پیش‌بینی قیمت برق پیشنهاد می‌کنند و آن را پرکردن یک شکاف علمی توصیف
می‌کنند: با وجود مدل‌های پیش‌بین متعدد، حوزه الگوریتم‌های یادگیری عمیق تا
آن زمان بررسی‌نشده مانده بود. نکته روش‌شناختی مهم در کار آنان، مقایسه
تجربی با الگوریتم‌های سنتی است \lr{---} یعنی همان رویه‌ای که سه سال بعد در
\cite{lago_forecasting_2021} به یک پروتکل رسمی تبدیل شد.

اوگورلو و همکاران \cite{ugurlu_electricity_2018} شبکه‌های عصبی بازگشتی را
به کار می‌گیرند و در مقدمه خود مشاهده‌ای را ثبت می‌کنند که برای این
پایان‌نامه اهمیت مستقیم دارد: به دلیل ماهیت چالش‌برانگیز قیمت برق
\lr{---} نوسان‌پذیری بالا، جهش‌های تیز و فصلی‌بودن \lr{---} انواع گوناگون
مدل‌ها همچنان با یکدیگر رقابت می‌کنند و \emph{هیچ‌یک نمی‌تواند به‌طور
مداوم بر دیگری برتری یابد}. این جمله، به‌تنهایی، توجیه‌کننده انتخاب
روش‌شناختی پژوهش حاضر است: اگر برتری مداوم وجود ندارد، آنگاه مقایسه بدون
آزمون معناداری آماری، چیزی جز نویز گزارش نمی‌کند.

ژانگ و همکاران \cite{zhang_forecasting_2018} رویکردی پیمانه‌ای برای
پیش‌بینی روز-پیش ارائه می‌دهند که بر داده تاریخی قیمت، داده جوی شبکه و
پیش‌بینی بار تکیه می‌کند و مطالعه موردی آن، بازار آلمان است. ترکیب
متغیرهای ورودی در این کار \lr{---} قیمت تاریخی به‌علاوه پیش‌بینی بار و
متغیرهای جوی \lr{---} به مجموعه ویژگی‌های پژوهش حاضر نزدیک است.

کوئو و همکاران \cite{kuo_electricity_2018} شبکه‌های عصبی عمیق با ساختار
ترکیبی را پیشنهاد می‌کنند و انگیزه خود را از منظر بازیگر بازار بیان
می‌کنند: راهبرد پیشنهاد قیمت، سطح سود تولیدکننده را تعیین می‌کند.

میلتیچ و همکاران \cite{miletic_day-ahead_2022} از شبکه‌های \lr{LSTM} برای
پیش‌بینی قیمت روز-پیش استفاده می‌کنند و پیامد عملی خطای پیش‌بینی را
صریح می‌سازند: پیش‌بینی نادرست به زمان‌بندی نادرست و کاهش سود می‌انجامد.
آنان افزایش سهم منابع تجدیدپذیر متغیر را عامل افزایش نوسان‌پذیری و
عدم‌قطعیت قیمت می‌دانند.

\subsection*{جمع‌بندی: قدرت بیشتر، تضمین کمتر}

الگوی مشترک این مطالعات آن است که مدل‌های عمیق‌تر، ظرفیت بیشتری برای
مدل‌سازی روابط غیرخطی دارند و در بسیاری از موارد دقت بالاتری گزارش
می‌کنند. اما دو ملاحظه در این خانواده برجسته‌تر از خانواده‌های پیشین است.

نخست، حساسیت به شرایط ارزیابی: هشدار اوگورلو و همکاران درباره نبود برتری
مداوم، به این معناست که رتبه‌بندی گزارش‌شده در یک مطالعه، لزوماً در مطالعه
دیگر بازتولید نمی‌شود. دوم، تفسیرپذیری: افزایش ظرفیت مدل با کاهش شفافیت
آن همراه است، که موضوع بخش \lr{2-6} است.

پژوهش حاضر هر دو ملاحظه را وارد طراحی می‌کند: ملاحظه نخست از طریق آزمون
معناداری آماری و آزمون برون‌توزیع، و ملاحظه دوم از طریق تحلیل
تفسیرپذیری مبتنی بر \lr{SHAP}.
